\documentclass[11pt]{article}
\usepackage[margin=1in]{geometry}
\usepackage{amsmath,amssymb,amsthm}
\usepackage{graphicx}
\usepackage{microtype}
\usepackage[hidelinks]{hyperref}

\newtheorem{theorem}{Theorem}
\newtheorem{corollary}[theorem]{Corollary}
\newtheorem{lemma}[theorem]{Lemma}
\newcommand{\R}{\mathbb{R}}

\title{Regulated Radial $L^p$ Multiplier Profiles Have Continuous Representatives}
\author{Candidate partial result for arXiv:2602.14623}
\date{11 August 2026}

\begin{document}
\maketitle

\noindent\textbf{Status.} Candidate partial result, likely valid. This note
does not solve the full radial multiplier continuity question.

\section{Source question}

Mikael de la Salle asks the following in Question 1.6 of
\cite{delasalle2026}, page 4:

\begin{quote}
Is there a discontinuous function $m:(0,\infty)\to\R$ such that the Fourier
multiplier with symbol $\xi\in\R^d\mapsto m(|\xi|)$ is bounded on
$L_p(\R^d)$ for some $p\ne 2$?
\end{quote}

Figure~\ref{fig:question} is a crop of the original source page.

\begin{figure}[ht]
  \centering
  \includegraphics[width=\textwidth]{figures/open_problem_crop.png}
  \caption{Question 1.6 in \cite[p.~4]{delasalle2026}.}
  \label{fig:question}
\end{figure}

The source proves that if the multiplier is bounded, then the primitive of
$t\mapsto m(e^t)$ is smooth in Zygmund's sense
\cite[Theorem~1.7]{delasalle2026}. Proposition 2.4 of the same paper gives the
equivalent weak-* tangent formulation used below.

\section{Definitions}

A function $u:I\to\R$ on an open interval is \emph{locally regulated} if at
every $t\in I$ the finite one-sided limits $u(t-)$ and $u(t+)$ exist. We do
not require either one-sided limit to equal the displayed point value $u(t)$.

Two bounded symbols represent the same Fourier multiplier when they agree
almost everywhere. Thus the appropriate continuity conclusion is the
existence of a continuous representative of the multiplier symbol.

\section{Idea of the proof}

The source theorem excludes every nonconstant weak-* tangent of the
logarithmic radial profile. A regulated jump has a canonical tangent: zooming
at the jump converges to the step function whose two values are the left and
right limits. The source theorem therefore forces those limits to agree at
every point. What remains are only removable point-value defects. A compact
interval can contain only finitely many defects of any fixed size, so their
union is countable and replacing them by the common two-sided limits produces
a continuous representative.

The highly oscillatory symbol singled out in Question 1.8 is not regulated
near the unit sphere, so this mechanism deliberately does not address it.

\section{Partial theorem and proof}

\begin{theorem}\label{thm:main}
Let $d\ge 2$, $1<p<\infty$, and $p\ne2$. Let $m\in L_\infty(0,\infty;\R)$
and assume that
\[
  \widehat{T_m f}(\xi)=m(|\xi|)\widehat f(\xi)
\]
defines a bounded operator on $L_p(\R^d)$. If the equivalence class of $m$
contains a locally regulated representative, then it contains a continuous
representative on $(0,\infty)$.
\end{theorem}

\begin{proof}
Choose a locally regulated representative of $m$. Composing it with the
continuous clipping map to an essential $L_\infty$ bound changes it only on a
null set and preserves regulatedness. We may therefore take this
representative to be bounded; still denote it by $m$, and put
\[
  u(t)=m(e^t),\qquad t\in\R.
\]
Then $u$ is bounded and locally regulated. By Theorem 1.7 and Proposition 2.4
of \cite{delasalle2026}, every weak-* limit in $\sigma(L_\infty,L_1)$ of a
sequence
\[
  x\longmapsto u(s_n+x/r_n),\qquad r_n\longrightarrow\infty,
\]
is constant almost everywhere.

Fix $s\in\R$ and set $a=u(s-)$ and $b=u(s+)$. Take $s_n=s$ and $r_n=n$.
For each fixed $x<0$,
\[
  u(s+x/n)\longrightarrow a,
\]
and for each fixed $x>0$ the same expression converges to $b$. The functions
are uniformly bounded, so dominated convergence against an arbitrary
$L_1(\R)$ test function shows weak-* convergence to
\[
  v(x)=a\mathbf 1_{(-\infty,0)}(x)+b\mathbf 1_{(0,\infty)}(x).
\]
The tangent characterization forces $v$ to be constant almost everywhere;
hence $a=b$. Since $s$ was arbitrary, the two one-sided limits of $u$ agree at
every point. Denote their common value by $c(s)$.

We next verify that $c$ is continuous. Given $s$ and $\varepsilon>0$, the
one-sided limit definitions give $\delta>0$ such that
\[
  |u(t)-c(s)|<\varepsilon
  \quad\text{whenever}\quad 0<|t-s|<\delta.
\]
If $s<t<s+\delta$, then points approaching $t$ from either side may be chosen
inside $(s,s+\delta)$, and passing to the corresponding one-sided limit gives
$|c(t)-c(s)|\le\varepsilon$. The same argument applies on the left. Thus $c$
is continuous.

It remains to show that $u=c$ almost everywhere. For integers $N,k\ge1$, set
\[
  E_{N,k}=\{s\in[-N,N]: |u(s)-c(s)|\ge 1/k\}.
\]
This set is finite. Indeed, if it had an accumulation point $s_0$, choose
distinct $s_j\in E_{N,k}$ with $s_j\to s_0$ and pass to a one-sided
subsequence. Regulatedness gives $u(s_j)\to c(s_0)$, while continuity gives
$c(s_j)\to c(s_0)$, contradicting
$|u(s_j)-c(s_j)|\ge1/k$. Therefore
$\{u\ne c\}=\bigcup_{N,k}E_{N,k}$ is countable, and $u=c$ almost everywhere.

Finally define $\widetilde m(r)=c(\log r)$. This function is continuous on
$(0,\infty)$ and equals $m(r)$ almost everywhere, so it represents the same
Fourier multiplier.
\end{proof}

\begin{corollary}
Under the multiplier assumptions of Theorem~\ref{thm:main}, the conclusion
holds whenever $m$ is locally of bounded variation, locally of finite
$q$-variation for some $1\le q<\infty$, or locally cadlag. In the cadlag case,
$m$ itself (rather than merely an almost-everywhere representative) is
continuous.
\end{corollary}

\begin{proof}
Each listed hypothesis implies local regulatedness. For finite
$q$-variation, failure of a one-sided limit would allow two interlaced
sequences with values separated by a fixed positive amount, forcing the sum
of the $q$th powers of successive increments to be infinite. Theorem
\ref{thm:main} applies. A cadlag function already equals its right-hand limit;
the proof of the theorem forces the left-hand limit to have the same value,
so the function is continuous at every point.
\end{proof}

\section{Verification notes}

\begin{itemize}
  \item The proof uses no unproved lemma. Its only external dependency is the
  weak-* tangent characterization in \cite[Theorem~1.7 and Proposition~2.4]
  {delasalle2026}.
  \item The step-function tangent is obtained in the full
  $\sigma(L_\infty,L_1)$ topology, not merely pointwise locally: dominated
  convergence applies to every $L_1(\R)$ test function.
  \item The representative issue is explicit. The proof does not claim that
  arbitrary null-set changes of a multiplier symbol are continuous.
  \item No computation is used or needed.
\end{itemize}

\noindent\textbf{Verifier verdict:} likely valid. The most useful human check
is the elementary compactness argument proving that $u=c$ almost everywhere.

\section{Scope and novelty status}

The full Question 1.6 remains open for bounded measurable profiles, and the
specific slowly oscillating profile in Question 1.8 is outside the regulated
class. The result is therefore only a regularity-class subcase.

On 11 August 2026, the local run indexes and parsed arXiv corpus were searched
for the exact arXiv id and close terminology. Bounded arXiv/web searches used
the exact id, the exact slowly oscillating candidate, and combinations of
``regulated'', ``one-sided limits'', ``bounded variation'', and ``radial
Fourier multiplier''. No explicit prior statement of this corollary was found.
Novelty confidence is modest because the proof is an elementary deduction
from Proposition 2.4 of the source paper itself.

\section{Human review recommendation}

Check the regulated-function lemma and retain the packet, if accepted, as a
narrow partial result. It should not be described as resolving the general
continuity question or Question 1.8.

\begin{thebibliography}{9}
\bibitem{delasalle2026}
Mikael de la Salle,
\emph{Kakeya conjecture and High-Rank Lattice von Neumann algebras},
arXiv:2602.14623v1, 2026.
\url{https://arxiv.org/abs/2602.14623}
\end{thebibliography}

\end{document}
